\section{Volume Hessians, finite elements, and computational complexity}
\label{sec:ipolybridge}

The variational description of a polyhedron has a second interpretation:
its infinitesimal deformation equation can be a discrete Dirichlet
problem. This is an exact connection to finite elements and electrical
networks, and it supplies the computational counterpart of the volume
and symplectic constructions developed above. The companion study
\emph{Duffin's resistor networks, the volume of ideal polyhedra, and
nearly-linear-time Laplacian solvers}~\cite{IpolyStudy} develops this
connection and its numerical examples. We incorporate the essential
identities here and connect them to boundary response, Gaussian
integration, flips, and the Casson--Rivin program.

\subsection{The cotangent form and the correct deformation coordinates}
\label{subsec:ipolyCoordinates}

A finite element calculation replaces a function by its values at mesh
vertices, extending it linearly across each triangle. Its energy is
then a quadratic form in finitely many numbers. An electrical network
has the same form: a conductance $w$ on an edge with voltage difference
$v_i-v_j$ contributes $w(v_i-v_j)^2$. Our task is to identify precisely
which polyhedral deformation variables produce this quadratic form.
The Hessian is the matrix of second derivatives; its negative measures
how rapidly volume decreases away from a maximum.

Let $\mathcal T$ be a finite, connected, consistently oriented simplicial
triangulation of a closed disk, with interior vertices $I$ and nonempty
boundary vertex set $B$. Each triangle has three distinct vertices.
For a piecewise affine function $v$ on a
geometric realization, integration over each triangle gives
\begin{equation}
 \int_{|\mathcal T|}|\nabla v|^2\,dA
 =\sum_{e=ij}w_e(v_i-v_j)^2=v^TLv,
 \qquad
 w_e=\frac12(\cot\alpha_e+\cot\beta_e).
 \label{eq:ipolyFEM}
\end{equation}
Here each undirected edge is counted once, $\alpha_e,\beta_e$ are its
opposite angles, and the second term is omitted on a boundary edge.
Here $L_{ij}=-w_{ij}$ for distinct joined vertices and
$L_{ii}=\sum_{j\ne i}w_{ij}$; missing edges have weight zero.
A barycentric function $\phi_i$ is the affine function equal to one
at vertex $i$ and zero at the other two vertices of a triangle.
The collection of these functions lets us reconstruct the piecewise
linear function from its vertex values.

To check the normalization, the
off-diagonal stiffness entry on a triangle is
$\int_T\nabla\phi_i\cdot\nabla\phi_j\,dA=-\tfrac12\cot\alpha_k$
for its barycentric functions. Row sums vanish, giving the displayed
edge form. This is the cotangent network of the finite element and
electrical interpretations~\cite{DuffinNetworks}.

The energy is nonnegative on every geometric mesh, even if some edge
weights are negative. The additional condition for nonnegative
\emph{interior} conductances is local Delaunayness:
\begin{equation}
 w_e=\frac{\sin(\alpha_e+\beta_e)}
              {2\sin\alpha_e\sin\beta_e}\ge0
 \quad\Longleftrightarrow\quad \alpha_e+\beta_e\le\pi.
 \label{eq:ipolyDelaunayweight}
\end{equation}
Grounding the boundary makes $L_{II}$ positive definite; a trianglewise
proof, including zero Delaunay weights, is given below. When the relevant
weights are nonnegative, solving its Dirichlet problem is solving
Kirchhoff's equations. Boundary-to-boundary edges do not enter
this Dirichlet system; positivity of their weights is an additional
condition if the full matrix is used in a different boundary problem.

Now prescribe the sums $\Delta_e=\alpha_e+\beta_e\le\pi$ opposite
interior edges, and prescribe the single opposite angle on each
boundary edge. An angle structure assigns positive corner angles,
with sum $\pi$ in every triangle and these prescribed edge sums.
Assume this positive angle domain is nonempty. For planar realization
the data also satisfy
$\sum_{e\ni i}(\pi-\Delta_e)=2\pi$ at each interior vertex; the
left-hand side is its total angle in every such angle structure.
If the intended planar boundary is strictly convex, impose also
\[
 0<\kappa_b:=\sum_{t\ni b}\alpha_{t,b}
   =\sum_{e\ni b}(\pi-\Delta_e)-\pi<\pi
 \qquad(b\in B).
\]
Here $\Delta_e$ denotes the single prescribed opposite angle on a
boundary edge. These boundary inequalities are needed for the convex
polyhedral interpretation; the Hessian identities below already hold
on the abstract positive angle domain. Positive feasibility alone
does not impose convexity of the developed boundary.
Orient the triangles consistently. If a corner $c$ is opposite the
directed edge $i\to j$, set
\begin{equation}
 \alpha_c(u)=\alpha_c^0+u_i-u_j,\qquad u|_B=0.
 \label{eq:ipolyAngleCoordinates}
\end{equation}
These are leading--trailing coordinates, not logarithmic scale factors.
Triangle sums telescope; increments opposite a shared edge cancel.
To prove that they parametrize the entire affine solution space, first
use the edge equations to leave one signed increment per interior edge.
The triangle equations are then the incidence equations of the connected
dual graph of faces, of rank $N_T-1$. If $E_{\rm int}$ is the number of
interior edges and $N_T$ the number of triangles, Euler's formula gives
\[
 \dim T\mathcal A=E_{\rm int}-N_T+1=|I|.
\]
The coordinate map is injective: if every corner increment vanishes,
$u$ is equal across every edge, hence constant by connectedness, and
the boundary condition makes it zero. It therefore has the required
dimension. This argument uses the disk topology; on an annulus extra
period directions can remain. In these coordinates the positive domain is
\[
 \Omega=\{u\in\mathbb R^I:\alpha_c^0+u_i-u_j>0
                  \text{ for every corner }c\}.
\]
It is nonempty, open, convex, and bounded. Positivity and the triangle
sums imply $\alpha_c<\pi$; injectivity transfers this boundedness to $u$.

Let $\Lambda'(a)=-\log(2\sin a)$ be the Lobachevsky derivative and put
$\mathcal V(\alpha)=\sum_c\Lambda(\alpha_c)$. At a geometric
configuration this is the sum of the volumes of the ideal tetrahedra
obtained by adjoining the vertex at infinity to the planar triangles.
Rivin's variational principle uses this concave functional to impose
geometric closure~\cite{RivinEuclidean,FuterGueritaud}. Existence and
realization statements require their geometric hypotheses; prescribed
Delaunay data and flat convex disks are treated separately in
\cite[Theorems~1.6 and~4.11]{RivinOptimization}. The differential
identities below do not assume that every point of $\Omega$ already
has a global planar realization.

\begin{proposition}[The exact finite-element Hessian]
\label{prop:ipolyHessian}
On the positive angle domain in the coordinates
\eqref{eq:ipolyAngleCoordinates},
\begin{equation}
 \nabla_u\mathcal V=(\log H_i)_{i\in I},\qquad
 -D_u^2\mathcal V=2L_{II}(\alpha),
 \label{eq:ipolyHessian}
\end{equation}
where $H_i$ is the dilation accumulated by developing the incident
triangles counterclockwise once around $i$, taking each radial edge to
the next. The matrix $L_{II}$ is positive definite when $I$ is nonempty.
In particular, at a geometric critical point
the negative volume Hessian is twice the finite element stiffness matrix.
\end{proposition}

\begin{proof}
In a positively oriented triangle $ijk$, differentiation with respect
to $u_i$ gives
$\log(\sin\alpha_{t,j}/\sin\alpha_{t,k})$.
In counterclockwise development from $ij$ to $ik$, the sine law gives
the next radial length divided by the previous one as
$\sin\alpha_{t,j}/\sin\alpha_{t,k}$. Summing the logarithms over the
star gives $\log H_i$. Reversing the traversal inverts $H_i$.
Since $\Lambda''(a)=-\cot a$, the two corners opposite an interior
edge contribute
$-(\cot\alpha_e+\cot\beta_e)(du_i-du_j)^2$ to $D^2\mathcal V$.
This is $-2$ times the corresponding summand of
\eqref{eq:ipolyFEM}. Boundary increments vanish.

For strictness, put $c_i=\cot\alpha_i$ on one positive triangle.
On $a_1+a_2+a_3=0$ its cotangent form is
\[
 \sum_i c_i a_i^2
  =(c_1+c_3)a_1^2+2c_3a_1a_2+(c_2+c_3)a_2^2.
\]
The first coefficient is
$c_1+c_3=\sin\alpha_2/(\sin\alpha_1\sin\alpha_3)>0$, and the
determinant is $c_1c_2+c_2c_3+c_3c_1=1$. Thus this form is positive
definite on the triangle tangent space, even if one cotangent is
negative. Applied to the three vertex differences, it vanishes exactly
on constant vertex values. Summing over the connected triangulation
shows that $\ker L$ consists of constants. Grounding $B$ removes this
kernel and proves the assertion. If $I$ is empty, the determinant of
the empty grounded matrix is understood to be one.
\end{proof}

Consequently $\mathcal V$ is strictly concave on $\Omega$. An interior
critical point, if it exists, is the unique maximum there and also a
maximum on its closure by continuity. Its equations $H_i=1$ kill
dilation holonomy. The separately imposed total-angle equations kill
rotational holonomy; convex boundary data and a global realization
argument supply the polyhedral interpretation. Strict concavity alone
does not place the maximizer in the interior.

For a true interior edge with $0<\Delta_e<\pi$, the weights satisfy
\begin{equation}
 w_e=\frac{\sin\Delta_e}{2\sin\alpha_e\sin\beta_e}
       \ge\cot(\Delta_e/2)>0.
 \label{eq:ipolyWeightLower}
\end{equation}
Thus positivity holds throughout the angle domain, before the angles
close to form a geometric mesh. At $\Delta_e=\pi$ the weight is zero;
at a geometric polyhedral realization such a diagonal subdivides a
flat face. A zero abstract weight by itself asserts no embedding.
The kernel proof also shows that these zero weights create no extra
grounded zero modes. The finite element interpretation
as an actual planar mesh applies at realization, whereas the Laplacian
identity and positive-network interpretation already apply to the
intermediate angle structures. Newton's equation is precisely
\begin{equation}
 2L_{II}(\alpha(u))\,\delta u=\log H(u).
 \label{eq:ipolyNewton}
\end{equation}
Figure~\ref{fig:mono-cotangent-hessian} summarizes the local edge weight
and its finite-element and volume-Hessian interpretations.

% BEGIN INLINE FIGURE: 03_cotangent_hessian.tex
\begin{figure}[htbp]
\centering
\begin{tikzpicture}[font=\small,line cap=round,>=Latex]
  \begin{scope}[shift={(-5.3,0)}]
    \coordinate (fi) at (0,0); \coordinate (fj) at (4,0);
    \coordinate (fk) at (1.5,2.6); \coordinate (fl) at (2.4,-2.4);
    \coordinate (fc1) at (2,.5788461538); \coordinate (fc2) at (2,-.4);
    \draw[mpblue!15] (fc1) circle (2.082081);\draw[mpblue!15] (fc2) circle (2.039608);
    \fill[mpblue!5] (fi)--(fj)--(fk)--cycle;
    \fill[mpteal!5] (fi)--(fl)--(fj)--cycle;
    \draw[mpgray!85] (fi)--(fk)--(fj) (fi)--(fl)--(fj);
    \draw[mpblue,line width=1.3pt] (fi)--(fj);
    \draw[mpochre,line width=1.7pt] (fc1)--(fc2);
    \draw[mpgray!60] (2.12,0)--(2.12,.12)--(2,.12);
    \pic[draw=mpgray,"$\alpha_e$",angle radius=7mm,angle eccentricity=1.5]
      {angle=fi--fk--fj};
    \pic[draw=mpgray,"$\beta_e$",angle radius=7mm,angle eccentricity=1.5]
      {angle=fj--fl--fi};
    \foreach \p in {fi,fj,fk,fl}\fill[mpblue] (\p) circle (1.7pt);
    \fill[mpochre] (fc1) circle (1.7pt) node[right=3pt] {$c_+$};
    \fill[mpochre] (fc2) circle (1.7pt) node[right=3pt] {$c_-$};
    \node[left] at (fi) {$i$};\node[right] at (fj) {$j$};
    \node[above] at (fk) {$k$};\node[below] at (fl) {$l$};
    \node[below,text=mpblue] at (1,0) {$e$};
    \node[left,text=mpochre] at (2,.32) {$e^*$};
  \end{scope}
  \node[align=left,text width=6.3cm,anchor=west] at (.35,.55) {%
    \textbf{One local weight, several roles}\\[7pt]
    $\displaystyle w_e=\frac12(\cot\alpha_e+\cot\beta_e)$\\[6pt]
    $\displaystyle\hphantom{w_e}=\frac{|e^*|_\pm}{|e|}$\\[9pt]
    $\displaystyle\int |\nabla u|^2=\sum_{e=ij}w_e(u_i-u_j)^2$\\[9pt]
    $\displaystyle -D_u^2\mathcal V=2L_{II}$\\[8pt]
    \textcolor{mpgray}{The last identity uses the specified angle
    coordinates and fixed boundary values.}};
\end{tikzpicture}
\caption[The cotangent Hessian and its finite-element interpretation]{The cotangent weight for two adjacent Euclidean triangles.
The dual segment joins their circumcenters and has signed length
$|e^*|_\pm$; in the locally Delaunay configuration shown, its sign is
positive. The finite-element stiffness matrix is the same cotangent
matrix that appears as minus one half of the volume Hessian in the
coordinates $\alpha_c(u)=\alpha_c^0+u_i-u_j$, with fixed edge-angle
sums and $u=0$ on the boundary. This identity is specific to that
coordinate and constraint system.}
\label{fig:mono-cotangent-hessian}
\end{figure}
% END INLINE FIGURE

\subsection{Boundary response, Schur complements, and a Gaussian factor}
\label{subsec:ipolyResponse}

The subscripts $I$ and $B$ mean interior and boundary vertices;
$L_{IB}$, for example, has interior rows and boundary columns.
A Dirichlet problem fixes boundary voltages and asks for interior
voltages with no injected current. The Neumann data are the currents
needed at the boundary to maintain those voltages.

For a three-vertex path with unit conductances, fix endpoint voltages
$0$ and $1$, and call the middle voltage $x$. The energy is
$x^2+(x-1)^2=2(x-1/2)^2+1/2$. Minimization gives $x=1/2$;
the effective conductance between the endpoints is $1/2$.
Completing this square is the one-variable version of the matrix
elimination below. Integrating $\exp[-(x^2+(x-1)^2)/(2\hbar)]$
over $x$ instead gives $\sqrt{\pi\hbar}\,e^{-1/(4\hbar)}$:
the same effective energy, now with a Gaussian prefactor.

For the finite element network, eliminating the interior voltage gives
\begin{equation}
 v_I=-L_{II}^{-1}L_{IB}v_B,\qquad
 \mathcal N=L_{BB}-L_{BI}L_{II}^{-1}L_{IB}.
 \label{eq:ipolyDtN}
\end{equation}
The matrix $\mathcal N$ maps boundary voltages to boundary currents;
it is the discrete Dirichlet-to-Neumann operator. Indeed, minimizing
$v^TLv$ at fixed $v_B$ gives the first equation and the minimum
$v_B^T\mathcal N v_B$. It follows that $\mathcal N$ is positive
semidefinite and that its kernel consists exactly of boundary constants.
Grounding one boundary vertex, or working modulo constants, makes it
positive definite. This energy statement allows signed boundary weights;
the full passive resistor interpretation requires nonnegative weights.

The same elimination appears in volume response. For a smooth function
with a nondegenerate local maximum at the origin, write
\[
 \mathscr V(x)=V_0-\tfrac12x^TKx+O(|x|^3),\qquad K>0,
\]
and let $Px$ be selected linear observables. These are geometric boundary
data only when the model supplies that identification; they are not the
fixed coordinates $u|_B=0$. The implicit-function theorem gives the local
maximizing branch $x(j)=K^{-1}P^Tj+O(|j|^2)$. Thus perturbation by a
source $j$, or restriction to $Px=b$, gives the local expansions
\begin{align}
 \max_{x\text{ near }0}\{\mathscr V(x)+j^TPx\}
    &=V_0+\tfrac12j^T\chi j+O(|j|^3),
 &\chi&=PK^{-1}P^T,\notag\\
 \max_{\substack{Px=b\\x\text{ near }0}}\mathscr V(x)
    &=V_0-\tfrac12b^T\chi^{-1}b+O(|b|^3),
 \label{eq:ipolyBoundaryResponse}
\end{align}
where the second formula assumes $P$ has full row rank, so $\chi>0$.
The first is the response to a source; the second fixes its conjugate
variable. They are related by a Legendre transform. In the planar
problem $K=2L_{II}$; in the Dehn filling problem, the same algebra gives
\eqref{eq:cuspHessian} and \eqref{eq:cuspmatrixHessian}. This identifies
an exact common operation without identifying their different spaces
of boundary data.

More generally, split $x=(y,z)$ and eliminate $y$. The effective
quadratic form is
\begin{equation}
 K_{\rm eff}=K_{zz}-K_{zy}K_{yy}^{-1}K_{yz},\qquad
 \det K=\det K_{yy}\det K_{\rm eff}.
 \label{eq:ipolySchur}
\end{equation}
For $\hbar>0$, completing the square gives the exact real Gaussian identity
\begin{equation}
 \int_{\mathbb R^{\dim y}}
 e^{-\frac1{2\hbar}(y,z)^TK(y,z)}\,dy
 =\frac{(2\pi\hbar)^{\dim y/2}}{\sqrt{\det K_{yy}}}
       e^{-z^TK_{\rm eff}z/(2\hbar)}.
 \label{eq:ipolyGaussian}
\end{equation}
Thus elimination in the Dirichlet problem, reduction of the classical
boundary action, and its leading Gaussian integration use the same
Schur complement. An oscillatory or complex phase requires the specified
contour and determinant branch, as in the later quantum calculations.
The finite-dimensional identity does not supply those analytic data
or identify an entire quantum invariant by itself.

There is also a combinatorial formula for the planar determinant. Let
$\mathfrak F_B$ be the spanning forests of the weighted graph in which
each component contains exactly one boundary vertex. Then
\begin{equation}
 \boxed{\det(-D_u^2\mathcal V)
      =2^{|I|}\sum_{F\in\mathfrak F_B}\prod_{e\in F}w_e.}
 \label{eq:ipolyForestDeterminant}
\end{equation}
To prove it, write $L_{II}=C_I\operatorname{diag}(w_e)C_I^T$, where
$C_I$ is the incidence matrix with boundary rows removed. Cauchy--Binet
expands its determinant over $|I|$-edge subsets. An incidence minor is
$\pm1$ exactly for a forest rooted at $B$, and zero otherwise. This
rooted-forest formula expresses the local Gaussian determinant in
geometric conductances. It is a deduction from the Hessian identity
and finite-dimensional linear algebra. The incidence-minor identity
holds for real signed weights as well; in the Delaunay Dirichlet problem
its contributing weights are nonnegative. Zero-weight edges contribute
zero terms, and strictness of $L_{II}$ makes the total positive.

\subsection{Why the nonlinear iteration needs a barrier}
\label{subsec:ipolyBarrierComplexity}

A barrier is a function that tends to $+\infty$ at the boundary
of the admissible domain. Minimizing $F-t\mathcal V$ balances
staying safely inside against increasing volume; increasing $t$
gradually reduces the influence of the barrier. A convex function
$f$ is \emph{self-concordant} here if
$|D^3f[h,h,h]|\le2(D^2f[h,h])^{3/2}$ for every direction $h$.
This compares the change of the Hessian with the local scale set
by the Hessian itself, which is what controls Newton steps.
The norm $\|h\|_H=\sqrt{h^THh}$ and its dual
$\|g\|_{H^{-1}}=\sqrt{g^TH^{-1}g}$ will be used below.

A fast linear solve does not by itself bound the number of Newton
steps. In particular, $-\mathcal V$ is not uniformly self-concordant:
near a vanishing angle, its entropy-like behavior makes the normalized
third derivative unbounded. The companion study supplies a more useful
estimate involving the barrier
\begin{equation}
 F(\alpha)=-\sum_c\log\sin\alpha_c.
 \label{eq:ipolyBarrier}
\end{equation}
The coefficient $\csc^2\alpha_c$ occurs in both $D^2F$ and
$D^3\mathcal V$, but these are tensors of different order. The precise
relation needed for optimization is the following inequality.

\begin{proposition}[A barrier compatible with ideal volume]
\label{prop:ipolyCompatibility}
For every positive angle structure and tangent vector $h$ whose
coordinates sum to zero in each triangle,
\begin{equation}
 |D^3\mathcal V[h,h,h]|
 \le\sqrt3\,(-D^2\mathcal V[h,h])\sqrt{D^2F[h,h]}.
 \label{eq:ipolyCompatibility}
\end{equation}
Consequently $F-t\mathcal V$ is self-concordant with the standard
constant for every $t\ge0$. The barrier $F$ has parameter at most
the number $m$ of corner angles.
\end{proposition}

\begin{proof}
For one triangle set $s_i=\sin\alpha_i$ and choose the cyclic side
directions $\psi=(0,\alpha_1+\alpha_2,\pi+\alpha_2)$.
They satisfy $\sum_i s_i e^{i\psi_i}=0$. Every vector with
$h_1+h_2+h_3=0$ can therefore be written
$h_i=\lambda s_i\cos\varphi_i$, where $\varphi_i=\psi_i-\tau$.
The cotangent energy identity and a trigonometric expansion give
\[
 A:=-D^2\mathcal V[h,h]=\lambda^2s_1s_2s_3,\quad
 B:=D^2F[h,h]=\lambda^2\sum_i\cos^2\varphi_i,
\]
\[
 D^3\mathcal V[h,h,h]
   =\lambda^3s_1s_2s_3\cos(\varphi_1+\varphi_2+\varphi_3).
\]
For the last identity, use $\cos^3 x=(3\cos x+\cos3x)/4$ and the
vanishing side-vector sum. Expanding the cosine of the sum shows
$|\cos(\varphi_1+\varphi_2+\varphi_3)|
\le\sum_i|\cos\varphi_i|\le\sqrt{3\sum_i\cos^2\varphi_i}$.
This proves the inequality on one triangle, and summation proves it
globally since $\sum A_T\sqrt{B_T}\le(\sum A_T)\sqrt{\sum B_T}$.
Each $-\log\sin\alpha_c$ is self-concordant, and, with
$A=-tD^2\mathcal V[h,h]$ and $B=D^2F[h,h]$,
\[
 |D^3(F-t\mathcal V)[h,h,h]|
 \le2B^{3/2}+\sqrt3 A\sqrt B\le2(A+B)^{3/2}.
\]
Finally Cauchy--Schwarz gives $|DF[h]|\le\sqrt m\sqrt{D^2F[h,h]}$,
the required barrier-parameter bound. Restriction to additional affine
angle constraints preserves these inequalities.
\end{proof}

All derivatives and inverse local norms in the following argument are
taken on the intrinsic affine angle space, after linear dependencies
have been removed. The parameter bound $m$ belongs to $F$;
$F-t\mathcal V$ is a standard self-concordant objective, not a barrier
claimed to have that same parameter.

Write $\Phi_t=F-t\mathcal V$ and $H_t=D^2\Phi_t$. A nonempty positive
angle domain has compact closure in angle coordinates. Since $F$ tends
to infinity at its relative boundary while $\mathcal V$ is continuous
on the closure, $\Phi_t$ has a unique interior minimizer $x_t$.
The zero-dimensional case is vacuous. If $\mathcal V^*$ denotes the
maximum on the closed feasible domain, then
\[
 \mathcal V^*-\mathcal V(x_t)\le\frac m t.
\]
Indeed, concavity of sine gives
$0\le\sin\beta_c\le\sin x_c+\cos x_c(\beta_c-x_c)$, hence
$DF(x)[\beta-x]\le m$ even for a closed feasible $\beta$ with zero
angles. At $x_t$, use $DF=tD\mathcal V$ and concavity of
$\mathcal V$. This argument does not assume that its closed maximizer
is positive.

For a finite iterate $x$, define the Newton decrement
$\lambda_t(x)=\|D\Phi_t(x)\|_{H_t(x)^{-1}}$. If $\lambda=\lambda_t(x)<1$,
self-concordance gives
$r:=\|x_t-x\|_{H_t(x)}\le\lambda/(1-\lambda)$.
For completeness, put $d=x_t-x$. Along the segment the one-direction
Hessian bound is
$D^2\Phi_t(x+s d)[d,d]\ge r^2/(1+s r)^2$.
Integrating and using stationarity at $x_t$ gives
$-D\Phi_t(x)[d]\ge r^2/(1+r)$, whereas dual Cauchy--Schwarz bounds
the left side by $\lambda r$. Moreover, $H_t\succeq D^2F$ implies
\[
 t\|D\mathcal V(x)\|_{H_t^{-1}}\le\sqrt m+\lambda.
\]
Concavity at $x$ now proves the finite-centering certificate
\[
 \mathcal V^*-\mathcal V(x)
 \le\frac{m+(\sqrt m+\lambda)\lambda/(1-\lambda)}{t}
 \le\frac{17m}{12t}\quad\text{if }\lambda\le\tfrac14,\ m\ge1.
\]

Here is an explicit short-step schedule. Start at $t_0>0$ with
$\lambda_{t_0}(x_0)\le1/4$, increase $t$ by
$1+1/(20\sqrt m)$, and take one full Newton step for the new objective.
With $\delta=(t'-t)/t$, the identities
\[
 D\Phi_{t'}=(1+\delta)D\Phi_t-\delta DF,
 \qquad H_{t'}\succeq H_t
\]
bound the new decrement before the step by
$(1+\delta)/4+\delta\sqrt m\le5/16$.
The standard self-concordant Newton estimate makes the step feasible
and bounds its new decrement by
$(\lambda/(1-\lambda))^2\le25/121<1/4$.
Stopping when $t\ge17m/(12\varepsilon)$ therefore gives
\begin{equation}
 O\!\left(\sqrt m\log\max\left\{e,\frac{m}{t_0\varepsilon}\right\}\right)
 \quad\hbox{Newton steps to volume gap at most }\varepsilon.
 \label{eq:ipolyStepBound}
\end{equation}
This bound assumes the stated initial centering; it does not pay for
finding or centering a feasible point. The underlying Newton estimates
are described in~\cite[Chapter~11]{BoydVandenberghe} and the authors'
lecture notes~\cite{BoydVandenbergheSlides}. They apply here to the
specified barrier schedule, not to unrestricted plain Newton iteration.

In the disk coordinates the Hessian of $F-t\mathcal V$ is again a
grounded planar Laplacian, now with interior conductances
\[
 2t w_e+\csc^2\alpha_e+\csc^2\beta_e.
\]
An energy-norm tolerance makes the inexact version concrete. At the new
parameter, let $d_N=-H^{-1}g$, $\lambda=\|d_N\|_H\le5/16$, and suppose
the computed step satisfies $\|d-d_N\|_H\le\lambda/100$.
Put $\eta=\|d-d_N\|_H$ and $r=\|d\|_H$. Then
$\eta\le1/320$ and $r\le101/320<1$. Hessian comparison and the
integrated gradient remainder give
\[
 \lambda_{\rm new}
 \le\frac{\eta}{1-r}+\frac{r^2}{(1-r)^2}
 \le\frac{10420}{47961}<\frac14.
\]
The old dual-norm linear residual is at most $\eta$, the nonlinear
remainder at most $r^2/(1-r)$, and conversion to the new dual norm
costs at most $1/(1-r)$. This argument assumes exact gradient and
Hessian evaluation apart from the stated solve error. Feasibility can
also be checked directly: $\|d\|_H<1$ and $H\succeq D^2F$ imply
$|d\alpha_c|<\sin\alpha_c\le\min\{\alpha_c,\pi-\alpha_c\}$ for
every corner, so the new angles remain strictly between zero and $\pi$.

Nearly-linear Laplacian solvers provide relative energy-norm
guarantees~\cite[Theorem~5.5]{SpielmanTengSolvers}. With such tolerances
maintained at every step, their combination with
\eqref{eq:ipolyStepBound} gives the conditional randomized expected
$\widetilde O(n^{3/2})$ arithmetic-work scale after centering, up to
accuracy factors. Here $n=|I|+|B|$ is the full mesh vertex count;
disk planarity gives $m=O(n)$ and $O(n)$ matrix entries.
Success probabilities over all solves,
coefficient precision and transcendental evaluations need their own
controls. This is not a bit-complexity or floating-point stability
bound. Finding a feasible point, centering it, and certifying the final
geometric error are separate costs. The companion's selected-input
timings do not bound these additional tasks. The dense implementation
described below follows the specified schedule, with numerical diagnostics;
it does not implement a nearly-linear solver or certify floating-point error.

\subsection{Flips connect the local matrix to mapping-class geometry}
\label{subsec:ipolyFlips}

A diagonal exchange has two different roles. Among triangulations of
the convex hull of a fixed planar vertex set, with fixed nodal values,
Rippa's theorem says that Delaunay
triangulation minimizes the piecewise-linear Dirichlet energy; in matrix
terms $L(\mathcal T_D)\preccurlyeq L(\mathcal T)$~\cite{RippaEnergy}.
This is an ordering of quadratic forms for the same geometric data.
It does not assert monotonicity of hyperbolic volume along an arbitrary
sequence of topological flips.
Figure~\ref{fig:mono-flip} illustrates the comparison for one fixed
Euclidean quadrilateral.

% BEGIN INLINE FIGURE: 08_flip.tex
\begin{figure}[htbp]
\centering
\begin{tikzpicture}[font=\small,>=Latex,line cap=round,line join=round]
  \begin{scope}[shift={(-5.3,0)}]
    \coordinate (fa) at (0,0);\coordinate (fb) at (3.2,0);
    \coordinate (fc) at (2.7,2);\coordinate (fd) at (.3,1.5);
    \draw[mpblue!15] (1.6,.6625) circle (1.731735);
    \fill[mpblue!5] (fa)--(fb)--(fc)--cycle;
    \fill[mpteal!5] (fa)--(fc)--(fd)--cycle;
    \draw[mpgray] (fa)--(fb)--(fc)--(fd)--cycle;
    \draw[mpochre,line width=1.35pt] (fa)--(fc);
    \foreach \p in {fa,fb,fc,fd}\fill[mpblue] (\p) circle (1.8pt);
    \node[below left] at (fa) {$A$};\node[below right] at (fb) {$B$};
    \node[above right] at (fc) {$C$};\node[above left] at (fd) {$D$};
    \node[font=\bfseries\small] at (1.6,2.85) {The diagonal $AC$};
    \node[text=mpochre,align=center,font=\scriptsize] at (1.6,-1.35)
      {Not locally Delaunay\\$D$ lies inside the circle through $A,B,C$};
  \end{scope}
  \draw[->,mpgray,line width=.85pt] (-1.1,.85)--(.9,.85)
    node[midway,above] {flip};
  \begin{scope}[shift={(1.8,0)}]
    \coordinate (ga) at (0,0);\coordinate (gb) at (3.2,0);
    \coordinate (gc) at (2.7,2);\coordinate (gd) at (.3,1.5);
    \fill[mpblue!5] (ga)--(gb)--(gd)--cycle;
    \fill[mpteal!5] (gb)--(gc)--(gd)--cycle;
    \draw[mpgray] (ga)--(gb)--(gc)--(gd)--cycle;
    \draw[mpteal,line width=1.35pt] (gb)--(gd);
    \foreach \p in {ga,gb,gc,gd}\fill[mpblue] (\p) circle (1.8pt);
    \node[below left] at (ga) {$A$};\node[below right] at (gb) {$B$};
    \node[above right] at (gc) {$C$};\node[above left] at (gd) {$D$};
    \node[font=\bfseries\small] at (1.6,2.85) {The diagonal $BD$};
    \node[text=mpteal,align=center,font=\scriptsize] at (1.6,-1.35)
      {Locally Delaunay\\the planar vertices have not moved};
  \end{scope}
  \node[align=center] at (0,-2.35)
    {For the same nodal values $u$: \quad
     $D_{\mathcal T_{BD}}(u)\le D_{\mathcal T_{AC}}(u)$.};
\end{tikzpicture}
\caption[A diagonal exchange and the Delaunay energy comparison]{A diagonal flip on a fixed Euclidean quadrilateral, shown
with its actual vertex positions. Rippa's energy comparison selects the
Delaunay diagonal for fixed nodal values. The same combinatorial move
also occurs in flip graphs used in surface topology; the displayed
energy comparison concerns the fixed Euclidean realization.}
\label{fig:mono-flip}
\end{figure}
% END INLINE FIGURE

For the next comparison, graph distance counts elementary moves.
A \emph{quasi-isometry} preserves large distances up to fixed
multiplicative and additive errors, with every target point within
a bounded distance of its image. It does not identify individual
edge lengths or local metrics. A mapping torus is formed from
$S\times[0,1]$ by identifying $(x,1)$ with $(f(x),0)$ for a surface
map $f$; its monodromy is the mapping class of $f$.
Pseudo-Anosov classes are those with the invariant stretching and
contracting measured foliations used in surface dynamics. The
comparison theorems below are external inputs, with their surface
and boundary hypotheses retained.

For a connected orientable finite-type marked surface admitting an ideal
triangulation, with at least one marked point on every boundary curve,
the graph of ideal triangulations and diagonal exchanges is instead
a coarse geometric model of its mapping class group. For orientation-
preserving classes and isotopies fixing the marked points individually,
the action is proper and cocompact, and an
orbit map is a quasi-isometry~\cite[Section~2.1, Lemma~2.3]{DisarloParlier}.
For a closed oriented surface of genus at least two, the pants graph is
a different model, quasi-isometric to Teichm\"uller space with its WP
metric. In this setting Brock compares the convex-core volume of the
quasi-Fuchsian manifold with conformal ends $X,Y$ to $d_{\rm WP}(X,Y)$,
with topology-dependent multiplicative and additive
constants~\cite[Theorems~1.1--1.2]{BrockWPVolumes}.
For pseudo-Anosov monodromy on the same closed surface, mapping-torus
volume and WP translation distance are comparable with topology-dependent
multiplicative constants~\cite[Theorem~1.1]{BrockMappingTori}.
The flip graph and pants graph have different metric roles.

There is a direct volume estimate when a flip sequence gives a layered
ideal triangulation of a hyperbolic mapping torus with $N$ tetrahedra:
straightening gives $\operatorname{Vol}(M)\le Nv_3$, where $v_3$ is the
volume of the regular ideal tetrahedron. This does not require the
given triangulation to have positive geometric shapes; signed
straightened volumes still have absolute value at most $v_3$.
It connects a combinatorial decomposition length to three-dimensional
volume. Comparing it further to WP geometry requires a theorem for the
same surface and boundary conventions, rather than transferring the
closed-surface theorem silently to a punctured surface.

\subsection{The three-manifold extension and the remaining questions}
\label{subsec:ipolyThreeManifolds}

The Casson--Rivin program uses the same ideal-tetrahedron volume and the
same angle-sum constraints~\cite{FuterGueritaud}. Its local Hessian is
still a cotangent quadratic form. For leading--trailing edge variables
$y_e$, put $p_\alpha=y_e+y_{e'}$ for the pair of opposite tetrahedron
edges carrying $\alpha$, and cyclically for $\beta,\gamma$. Then
\begin{equation}
 -D^2\mathcal V
 =\sum_t\left[
 \cot\alpha_t(dp_\gamma-dp_\beta)^2
 +\cot\beta_t(dp_\alpha-dp_\gamma)^2
 +\cot\gamma_t(dp_\beta-dp_\alpha)^2\right].
 \label{eq:ipolyThreeDimensionalHessian}
\end{equation}
This displays the Euclidean cusp-link triangle inside the second
variation. But a square now involves sums of opposite-edge variables,
and the assembled matrix need not have nonpositive off-diagonal entries
or be diagonally dominant. Additional peripheral directions and removal
of linear dependencies must also be included in a full coordinate
system. The planar positive-Laplacian solver conclusion therefore does
not transfer automatically to arbitrary three-manifold triangulations.
The trianglewise barrier estimate does transfer, because it survives
affine restriction.

For an ideal triangulation of an orientable compact three-manifold whose
nonempty boundary is a union of tori, a positive angle structure critical
for volume on the full angle space supplies the complete hyperbolic
metric on its interior, by the Casson--Rivin
theorem~\cite[Theorem~1.2]{FuterGueritaud}. The angle space has triangle
sums $\pi$ and edge sums $2\pi$; full criticality includes the peripheral
directions. Edge variations alone do not impose completeness.
A boundary maximizer need not realize the
chosen triangulation by positive ideal tetrahedra. Optimization on a
given angle polytope therefore does not solve the problem of finding
a suitable geometric triangulation or yield a polynomial-time
geometrization algorithm.

The finite element comparison also does not identify its real symmetric
Hessian with the alternating WP form. The exact links established here
are the volume Hessian, its inverse response, its Schur complement,
and its Gaussian determinant. A continuum limit relating appropriately
renormalized polyhedral functionals to a WP potential would require
geometric convergence, a specified normalization, and control of second
variations. This is a natural question, not a theorem asserted here.
Likewise, in three dimensions it remains useful to ask for an effective
bound on the smallest angle at an interior maximum in terms of the
number of tetrahedra and a quantitative feasible margin. Such a bound
would connect variational existence to numerical certification and
complexity. These questions describe how the computational strand can
extend the monograph's variational and semiclassical program.

\subsection{Reproducible computational examples}
\label{subsec:ipolyComputation}

The \href{../experiments/section9/README.md}{Section 9 reproduction guide}
connects these formulas to a fixed $3$ by $2$ integer-coordinate grid,
triangulated into twelve right-isosceles triangles with two interior
vertices. At its geometric angle assignment,
\[
 L_{II}=\begin{pmatrix}4&-1\\-1&4\end{pmatrix},\qquad
 K=-D_u^2\mathcal V=\begin{pmatrix}8&-2\\-2&8\end{pmatrix},\qquad
 \det K=60.
\]
An independent rational finite-element assembly gives this matrix.
Exhaustive boundary-rooted forest enumeration gives fifteen unit-weight
terms, so \eqref{eq:ipolyForestDeterminant} gives $2^2\cdot15=60$.
Eliminating the first variable of $K$ leaves $K_{\rm eff}=15/2$, and
$8(15/2)=60$. Prescribing the electrical boundary voltage $v=x/3$
gives interior values $1/3,2/3$ and energy $2/3$. These are exact finite
identities for the named fixture, separate from the general proofs above.

The \href{../packages/volume-toolkit/ideal_poly_volume_toolkit/variational.py}
{maintained numerical operators} assemble angle increments, gradients,
Hessians and boundary response. The
\href{../experiments/section9/run.py}{experiment driver} also compares
the binary64 volume with $12G$, where $G$ is Catalan's constant, and an
off-origin Hessian with independent numerical differentiation at 50-digit
working precision. The \href{../artifacts/section9/result.json}{result record}
states source and input hashes, methods, residuals and tolerances.
These numerical comparisons are not interval enclosures.

The \href{../site/section9/index.html}{worked computational companion}
and \href{../site/section9/downloads/section9-reproduction.zip}{complete
reproduction archive} provide the input, runnable code and environment
requirements. The grid has zero-weight square diagonals and is therefore
non-strict Delaunay. Off-origin angles belong to the abstract domain;
their geometric closure is not asserted.

The separate \href{../site/section9/barrier.html}{barrier examples} exercise
the \href{../packages/volume-toolkit/ideal_poly_volume_toolkit/angle_barrier.py}
{dense affine-angle solver}. It accepts a supplied positive angle vector
and independent tangent columns, performs separately budgeted initial
centering, and then follows the short-step parameter schedule above.
Every accepted point records its parameter, Newton decrement, angle margins,
linear-solve residual and numerical conditioning. Rejected domains, failed
steps and exhausted budgets remain unsuccessful outcomes; angles are not
clipped into the domain. The reported finite-gap estimate is evaluated in
binary64 arithmetic and is not an interval certificate.

The \href{../experiments/section9/barrier_run.py}{barrier experiment driver}
compares the computed endpoints with independent multiprecision volume
references and checks the recorded path and stopping conditions. The fixed
cases are the full single-triangle angle simplex with its maximum at the
regular ideal tetrahedron, an asymmetric rectangular disk, and a two-triangle
affine slice with one flat triangle at its closed maximum. Its
\href{../artifacts/section9/barrier-result.json}{result record} separates
initial-centering work from path steps, and its
\href{../site/section9/downloads/barrier-reproduction.zip}{reproduction archive}
contains the actual code, inputs and environment requirements. These finite
examples do not establish a scaling law. A nearly-linear solver, a certified
relative energy-error bound, a peripheral-holonomy solver and a global
realization certificate remain outside this implementation. The
\href{claims/section9.json}{claim ledger} distinguishes this numerical
coverage from the general proofs, external theorem inputs and open questions.

